\documentclass[11pt]{article}

\usepackage[T1]{fontenc}
\usepackage{lmodern}
\usepackage{geometry}
\usepackage{hyperref}
\usepackage{xcolor}
\usepackage{listings}
\usepackage{longtable}
\usepackage{booktabs}
\usepackage{multicol}
\usepackage{titlesec}

\geometry{margin=1in}

\definecolor{codebg}{RGB}{245,245,245}
\definecolor{keyword}{RGB}{0,0,150}

\lstset{
  basicstyle=\ttfamily\small,
  backgroundcolor=\color{codebg},
  frame=single,
  columns=fullflexible,
  keepspaces=true
}

\titleformat{\section}{\Large\bfseries}{}{0em}{}
\titleformat{\subsection}{\large\bfseries}{}{0em}{}

\title{Bash Parameter Expansion Cheat Sheet}
\author{}
\date{}

\begin{document}
\maketitle

\tableofcontents
\newpage

\section{Why Parameter Expansion Matters}

Bash parameter expansion performs string manipulation entirely inside the shell.
This avoids process creation overhead from utilities such as \texttt{basename},
\texttt{dirname}, \texttt{cut}, \texttt{sed}, and \texttt{awk}.

In tight loops or scripts processing tens of thousands of strings, replacing
external commands with parameter expansion can reduce runtime from minutes to seconds.

\section{General Notes}

\begin{itemize}
\item Operates on shell variables only.
\item Patterns use shell globs, not regular expressions.
\item Always quote expansions unless word splitting is desired.
\item Best suited for simple delimiter- and suffix-based transforms.
\end{itemize}

\section{Prefix and Suffix Removal}

\subsection{Syntax Summary}

\begin{longtable}{lll}
\toprule
Form & Meaning & Match Type \\
\midrule
\texttt{\$\{v\#pattern\}}   & Remove from start & Shortest \\
\texttt{\$\{v\#\#pattern\}} & Remove from start & Longest \\
\texttt{\$\{v\%pattern\}}   & Remove from end   & Shortest \\
\texttt{\$\{v\%\%pattern\}} & Remove from end   & Longest \\
\bottomrule
\end{longtable}

\subsection{Examples}

\begin{lstlisting}
v="aa/bb/cc.txt"

${v#*/}      -> bb/cc.txt
${v##*/}     -> cc.txt
${v%/*}      -> aa/bb
${v%%/*}     -> aa
\end{lstlisting}

\subsection{Comparison}

\begin{tabular}{ll}
\toprule
Parameter Expansion & External Utility \\
\midrule
\texttt{\$\{p\#\#*/\}} & \texttt{basename p} \\
\texttt{\$\{p\%/*\}} & \texttt{dirname p}  \\
\bottomrule
\end{tabular}

\section{Basename, Dirname, Extensions}

\subsection{Common Patterns}

\begin{lstlisting}
p="/var/log/nginx/access.log"

base=${p##*/}     # access.log
dir=${p%/*}       # /var/log/nginx
stem=${base%.*}   # access
ext=${base##*.}   # log
\end{lstlisting}

\subsection{Utility Comparison}

\begin{lstlisting}
basename -- "$p"
dirname  -- "$p"
\end{lstlisting}

Parameter expansion is dramatically faster when used repeatedly inside loops.

\section{Substring Extraction}

\subsection{By Position}

\begin{lstlisting}
s="abcdef"

${s:2}     -> cdef
${s:2:3}   -> cde
${s: -2}   -> ef
\end{lstlisting}

\subsection{Comparison}

\begin{lstlisting}
awk '{print substr($0,3,3)}'
\end{lstlisting}

Prefer expansion when positions are fixed and data is already in variables.

\section{Search and Replace}

\subsection{Syntax}

\begin{longtable}{ll}
\toprule
Form & Meaning \\
\midrule
\texttt{\$\{v/pat/repl\}}    & Replace first match \\
\texttt{\$\{v//pat/repl\}}   & Replace all matches \\
\texttt{\$\{v/\#pat/repl\}}  & Replace at start \\
\texttt{\$\{v/\%pat/repl\}}  & Replace at end \\
\bottomrule
\end{longtable}

\subsection{Examples}

\begin{lstlisting}
v="a-b-c"

${v/-/_}     -> a_b-c
${v//-/_}    -> a_b_c
\end{lstlisting}

\subsection{Comparison}

\begin{lstlisting}
sed 's/-/_/'
sed 's/-/_/g'
\end{lstlisting}

Use \texttt{sed} if regular expressions or capture groups are required.

\section{Case Conversion}

(Bash 4+)

\begin{lstlisting}
v="MiXeD"

${v,,}   -> mixed
${v^^}   -> MIXED
${v,}    -> miXeD
${v^}    -> MiXeD
\end{lstlisting}

\subsection{Comparison}

\begin{lstlisting}
tr '[:upper:]' '[:lower:]'
\end{lstlisting}

\section{Defaults, Errors, and Conditionals}

\begin{longtable}{ll}
\toprule
Form & Meaning \\
\midrule
\texttt{\$\{v:-d\}} & Use default if unset or empty \\
\texttt{\$\{v:=d\}} & Assign default if unset or empty \\
\texttt{\$\{v:?msg\}} & Error if unset or empty \\
\texttt{\$\{v:+alt\}} & Use alt if set \\
\bottomrule
\end{longtable}

\section{Length}

\begin{lstlisting}
len=${#v}
\end{lstlisting}

\subsection{Comparison}

\begin{lstlisting}
wc -c <<<"$v"
\end{lstlisting}

\section{Whitespace Trimming (Advanced)}

Pure parameter expansion trimming is possible but verbose:

\begin{lstlisting}
# leading
v="${v#"${v%%[!$' \t\r\n']*}"}"
# trailing
v="${v%"${v##*[!$' \t\r\n']}"}"
\end{lstlisting}

For readability, consider wrapping this in a function.

\section{When to Use What}

\begin{longtable}{ll}
\toprule
Use Parameter Expansion & Use External Tools \\
\midrule
High iteration loops & Complex regex parsing \\
Simple delimiters & Multi-field arithmetic \\
Path manipulation & Streaming input \\
Suffix/prefix logic & Readability over speed \\
\bottomrule
\end{longtable}

\section{Mental Model Summary}

\begin{itemize}
\item \texttt{\#} removes from the left
\item \texttt{\%} removes from the right
\item doubling (\texttt{\#\#}, \texttt{\%\%}) means ``greedy''
\item replacements are glob-based, not regex-based
\end{itemize}

\end{document}
